% =====================================================================
% Modulprojekt DLT – Bericht (ACL-Vorlage)
%
% SO BENUTZEN SIE DIESE DATEI:
% 1. Gehen Sie zu Overleaf -> New Project -> All Templates -> Suche nach "Association for Computational Linguistics" -> Auswahl "Association for Computational Linguistics (ACL) conference" -> Open As Template.
%    Die Vorlage enthaelt bereits acl.sty und acl_natbib.bst.
% 2. Ersetzen Sie die dortige Haupt-.tex durch DIESE Datei und legen Sie
%    custom.bib daneben.
% 3. Kompilieren mit pdfLaTeX. Richtwert: max. ~6 Seiten (ohne Literatur/Anhang).
%
% Alternativ lokal: die Style-Dateien liegen unter
%   github.com/acl-org/acl-style-files (latex/acl.sty, latex/acl_natbib.bst)
%
% WICHTIG – UMFANG DES BERICHTS:
%   Der Bericht behandelt AUSSCHLIESSLICH Teil B (Ihre eigene Übertragung auf
%   eine andere Aufgabenart). Teil A (das rotten_tomatoes-Notebook) wird NICHT
%   im Bericht beschrieben – Sie reichen dieses Notebook aber ausgeführt mit ein.
%   Ersetzen Sie unten überall den Platzhalter <IHR TEIL-B-DATENSATZ>.
% =====================================================================
\documentclass[11pt]{article}
\usepackage[review]{acl}   % Fuer die finale Abgabe: [review] entfernen.

\usepackage{times}
\usepackage{latexsym}
\usepackage[T1]{fontenc}
\usepackage[utf8]{inputenc}
\usepackage{microtype}
\usepackage{graphicx}
\usepackage{booktabs}

\title{Feintuning von BERT für <IHRE TEIL-B-AUFGABE> \\
       auf dem Datensatz <IHR TEIL-B-DATENSATZ>}

\author{Vorname Nachname \\
  Universität Potsdam \\
  \texttt{email@uni-potsdam.de} \\}

\begin{document}
\maketitle

\begin{abstract}
Kurze Zusammenfassung (4--6 Sätze): Aufgabe, Datensatz, Methode, wichtigstes
Ergebnis und eine Erkenntnis. Schreiben Sie den Abstract zuletzt.
\end{abstract}

% Dieser Bericht behandelt ausschliesslich Teil B (die Uebertragung auf eine
% andere Aufgabenart). Teil A wird nicht beschrieben, das Notebook aber
% ausgefuehrt mit eingereicht.

\section{Einleitung}
Motivieren Sie Ihre \emph{Teil-B-Aufgabe} (z.\,B.\ Paraphrasenerkennung oder
NER) und formulieren Sie in ein bis zwei Sätzen die Fragestellung Ihrer
Untersuchungen (Abschnitt~8 des Notebooks). Nennen Sie am Ende der Einleitung
kurz Ihr Hauptergebnis. Teil~A (rotten\_tomatoes) ist \textbf{nicht} Gegenstand
des Berichts.

\section{Verwandte Arbeiten}
Zwei bis drei Referenzen genügen: das Modell \citep{devlin2019bert} bzw.
\citep{sanh2019distilbert}, der Datensatz \citep{pang2005seeing} und ggf.
die Methode Ihrer Untersuchung, z.\,B.\ LoRA \citep{hu2022lora}.

\section{Daten}
Beschreiben Sie Ihren \emph{Teil-B-Datensatz}: Größe der Splits, Klassenbalance
(bzw.\ Label-/Entity-Verteilung), Textlängen. Verweisen Sie auf Ihre
Abbildung/Statistiken aus dem Notebook.

\section{Methode}
Modell (DistilBERT), Tokenisierung (\texttt{max\_length}), Klassifikationskopf,
Trainingsdetails (Lernrate, Epochen, Batch-Größe, Seed). Beschreiben Sie auch
Ihre beiden Baselines (Majority-Class, TF-IDF + logistische Regression).

\section{Experimente}
Aufbau: Welche Splits wofür? Welche Metriken (Accuracy, Macro-F1) und warum?
Beschreiben Sie das Setup Ihrer eigenen Untersuchung so, dass es
reproduzierbar ist.

\section{Ergebnisse}
Mindestens eine Tabelle (Tabelle~\ref{tab:results}) und eine Abbildung.
Vergleichen Sie BERT gegen \emph{beide} Baselines.

\begin{table}[t]
\centering
\begin{tabular}{lcc}
\toprule
Modell & Accuracy & Macro-F1 \\
\midrule
Majority-Class      & 0.00 & 0.00 \\
TF-IDF + LogReg      & 0.00 & 0.00 \\
DistilBERT (fein)    & 0.00 & 0.00 \\
\bottomrule
\end{tabular}
\caption{Ergebnisse auf dem Test-Set. Werte aus \texttt{results/metrics.json}.}
\label{tab:results}
\end{table}

\section{Fehleranalyse}
Diskutieren Sie einige falsch klassifizierte Beispiele und Ihre Hypothesen
(Negation, Ironie, Länge, \dots).

\section{Diskussion und Limitationen}
Was hat funktioniert, was nicht? Grenzen des Aufbaus (kleiner Datensatz, ein
Seed, kein Hyperparameter-Sweep). Ein Ausblick in ein bis zwei Sätzen.

\section*{Reproduzierbarkeit}
Link zum Code-Repository und die exakte Anweisung zum Reproduzieren
(Notebook/Skript, Seed).

% ---------------------------------------------------------------------
\bibliography{custom}

\end{document}
